% GENERATED FILE. Edit canonical lessons, then run npm run generate:books.
% canonical-source-hash: fnv1a64:75d6d3b0635a83e0
% canonical-lessons: MR-C173-tarihi, MR-C173-jari, MR-C173-asunhi, MR-C173-durdaivane, MR-C173-sudaivane

\chapter{Even so, Although, Despite, Unfortunately, Luckily}
\label{ch:mr-even-so-although-despite-unfortunately-luckily}
\hlchaptermodality{173}

\begin{chapteropening}
\textbf{By the end of this chapter:} \emph{I can say even so, although, despite, unfortunately and luckily.}

Complete the last lesson of chapter 173: I can say even so, although, despite, unfortunately and luckily.
\end{chapteropening}

\section[\texorpdfstring{\mr{तरीही} (tarīhī)}{तरीही (tarīhī)}]{\mr{तरीही} (tarīhī) — even so}

\label{lesson:MR-C173-tarihi}



\begin{quote}
Before the new one: say the Marathi for over there, then the Marathi for in the middle.
\end{quote}

\subsection*{You'll want to know: \mr{तरीही}}
\textbf{\mr{तरीही}} — \emph{tarīhī} — \textquotedblleft{}even so\textquotedblright{}.

\begin{cousinweb}
\textbf{\mr{तरीही}} (\emph{tarīhī}) is \emph{tarī}, which CDIAL lists, as \textquotedblleft{}then\textquotedblright{}, under Sanskrit \emph{tarhi}, \textquotedblleft{}then\textquotedblright{} (CDIAL 5730), with \emph{hī}, \textquotedblleft{}even\textquotedblright{}.
\end{cousinweb}

\begin{grammarlens}[title={What you've built}]
One more word for this chapter.
\end{grammarlens}

\subsection*{Your turn}
\begin{itemize}
\raggedright
  \item \emph{Say it:} \emph{tarīhī}
  \item \emph{Say it:} \emph{tarīhī}, once more
  \item \emph{Say it:} say \emph{tarīhī}
  \item \emph{Recall:} say the Marathi for over there, then the Marathi for in the middle, then say \emph{tarīhī} again
\end{itemize}

\begin{tcolorbox}[breakable,colback=teal!4,colframe=teal!35!black,title={Before you move on}]
What does \mr{तरीही} mean? (\textquotedblleft{}Even so\textquotedblright{}.) Say it once more.
\end{tcolorbox}
\section[\texorpdfstring{\mr{जरी} (jarī)}{जरी (jarī)}]{\mr{जरी} (jarī) — although}

\label{lesson:MR-C173-jari}



\begin{quote}
Before the new one: say the Marathi for in the middle, then the Marathi for even so.
\end{quote}

\subsection*{You'll want to know: \mr{जरी}}
\textbf{\mr{जरी}} — \emph{jarī} — \textquotedblleft{}although\textquotedblright{}.

\begin{cousinweb}
\textbf{\mr{जरी}} (\emph{jarī}) is a tadbhava, a word inherited from Sanskrit and reshaped by sound change. It comes from \emph{yarhi}, \textquotedblleft{}when\textquotedblright{} (CDIAL 10429). CDIAL glosses the Marathi word \textquotedblleft{}if\textquotedblright{}.
\end{cousinweb}

\begin{grammarlens}[title={What you've built}]
One more word for this chapter.
\end{grammarlens}

\subsection*{Your turn}
\begin{itemize}
\raggedright
  \item \emph{Say it:} \emph{jarī}
  \item \emph{Say it:} \emph{jarī}, once more
  \item \emph{Say it:} say \emph{jarī}
  \item \emph{Recall:} say the Marathi for in the middle, then the Marathi for even so, then say \emph{jarī} again
\end{itemize}

\begin{tcolorbox}[breakable,colback=teal!4,colframe=teal!35!black,title={Before you move on}]
What does \mr{जरी} mean? (\textquotedblleft{}Although\textquotedblright{}.) Say it once more.
\end{tcolorbox}
\section[\texorpdfstring{\mr{असूनही} (asūnhī)}{असूनही (asūnhī)}]{\mr{असूनही} (asūnhī) — despite, even though}

\label{lesson:MR-C173-asunhi}



\begin{quote}
Before the new one: say the Marathi for even so, then the Marathi for although.
\end{quote}

\subsection*{You'll want to know: \mr{असूनही}}
\textbf{\mr{असूनही}} — \emph{asūnhī} — \textquotedblleft{}despite, even though\textquotedblright{}.

\begin{cousinweb}
\textbf{\mr{असूनही}} (\emph{asūnhī}) is \emph{asūn}, \textquotedblleft{}being\textquotedblright{}, from the verb \emph{asṇẽ}, \textquotedblleft{}to be\textquotedblright{}, with \emph{hī}, \textquotedblleft{}even\textquotedblright{}. CDIAL lists \emph{asṇẽ} under Sanskrit \emph{āste}, \textquotedblleft{}sits, exists, remains\textquotedblright{} (CDIAL 1480), or perhaps \emph{ākṣeti}, \textquotedblleft{}abides\textquotedblright{} (CDIAL 1031).
\end{cousinweb}

\begin{grammarlens}[title={What you've built}]
One more word for this chapter.
\end{grammarlens}

\subsection*{Your turn}
\begin{itemize}
\raggedright
  \item \emph{Say it:} \emph{asūnhī}
  \item \emph{Say it:} \emph{asūnhī}, once more
  \item \emph{Say it:} say \emph{asūnhī}
  \item \emph{Recall:} say the Marathi for even so, then the Marathi for although, then say \emph{asūnhī} again
\end{itemize}

\begin{tcolorbox}[breakable,colback=teal!4,colframe=teal!35!black,title={Before you move on}]
What does \mr{असूनही} mean? (\textquotedblleft{}Despite, even though\textquotedblright{}.) Say it once more.
\end{tcolorbox}
\section[\texorpdfstring{\mr{दुर्दैवाने} (durdaivāne)}{दुर्दैवाने (durdaivāne)}]{\mr{दुर्दैवाने} (durdaivāne) — unfortunately}

\label{lesson:MR-C173-durdaivane}



\begin{quote}
Before the new one: say the Marathi for although, then the Marathi for despite.
\end{quote}

\subsection*{You'll want to know: \mr{दुर्दैवाने}}
\textbf{\mr{दुर्दैवाने}} — \emph{durdaivāne} — \textquotedblleft{}unfortunately\textquotedblright{}.

\begin{cousinweb}
\textbf{\mr{दुर्दैवाने}} (\emph{durdaivāne}) is Sanskrit \emph{durdaiva}, \textquotedblleft{}ill luck\textquotedblright{}: \emph{dur-}, \textquotedblleft{}bad\textquotedblright{}, and \emph{daiva}, \textquotedblleft{}fate\textquotedblright{}, with the ending \textbf{-\mr{ने}}, \textquotedblleft{}by\textquotedblright{}.
\end{cousinweb}

\begin{grammarlens}[title={What you've built}]
One more word for this chapter.
\end{grammarlens}

\subsection*{Your turn}
\begin{itemize}
\raggedright
  \item \emph{Say it:} \emph{durdaivāne}
  \item \emph{Say it:} \emph{durdaivāne}, once more
  \item \emph{Say it:} say \emph{durdaivāne}
  \item \emph{Recall:} say the Marathi for although, then the Marathi for despite, then say \emph{durdaivāne} again
\end{itemize}

\begin{tcolorbox}[breakable,colback=teal!4,colframe=teal!35!black,title={Before you move on}]
What does \mr{दुर्दैवाने} mean? (\textquotedblleft{}Unfortunately\textquotedblright{}.) Say it once more.
\end{tcolorbox}
\section[\texorpdfstring{\mr{सुदैवाने} (sudaivāne)}{सुदैवाने (sudaivāne)}]{\mr{सुदैवाने} (sudaivāne) — luckily}

\label{lesson:MR-C173-sudaivane}



\begin{quote}
Before the new one: say the Marathi for despite, then the Marathi for unfortunately.
\end{quote}

\subsection*{You'll want to know: \mr{सुदैवाने}}
\textbf{\mr{सुदैवाने}} — \emph{sudaivāne} — \textquotedblleft{}luckily\textquotedblright{}.

\begin{cousinweb}
\textbf{\mr{सुदैवाने}} (\emph{sudaivāne}) is Sanskrit \emph{sudaiva}, \textquotedblleft{}good luck\textquotedblright{}: \emph{su-}, \textquotedblleft{}good\textquotedblright{}, and \emph{daiva}, \textquotedblleft{}fate\textquotedblright{}, with the ending \textbf{-\mr{ने}}, \textquotedblleft{}by\textquotedblright{}.
\end{cousinweb}

\begin{grammarlens}[title={What you've built}]
One more word for this chapter.
\end{grammarlens}

\subsection*{Your turn}
\begin{itemize}
\raggedright
  \item \emph{Say it:} \emph{sudaivāne}
  \item \emph{Say it:} \emph{sudaivāne}, once more
  \item \emph{Say it:} say \emph{sudaivāne}
  \item \emph{Recall:} say the Marathi for despite, then the Marathi for unfortunately, then say \emph{sudaivāne} again
\end{itemize}

\begin{tcolorbox}[breakable,colback=teal!4,colframe=teal!35!black,title={Before you move on}]
What does \mr{सुदैवाने} mean? (\textquotedblleft{}Luckily\textquotedblright{}.) Say it once more.
\end{tcolorbox}
